\documentclass[12pt]{article}
\usepackage[utf8]{inputenc}
\usepackage[T1]{fontenc}
\usepackage{hyperref}
\title{Poqomchi' Language Resources for K-12 Curriculum Development}
\author{Compiled by Claude Sonnet 5.5 (claude-sonnet-5-5) for Andr\'as Kornai}
\date{}
\begin{document}
\maketitle

\section{Introduction}
This dossier covers Western Poqomchi' (ISO 639-3 pob; Glottocode poqo1253 is not assumed here, the identification given was poqo1256), a Mayan language of Guatemala (Alta Verapaz, Baja Verapaz and El Quich\'e). Poqomchi' as a whole is reported to have roughly 130,000 speakers. Many of the sources found do not say whether they describe the Western or the Eastern variety, so that is flagged for each item. The search was short and little pedagogical material turned up.

\section{Primary Educational Resources}
\begin{itemize}
\item Academia de las Lenguas Mayas de Guatemala (ALMG), bilingual Poqomchi'--Spanish \emph{Diccionario Pedag\'ogico} (reported May 2025), aimed at classroom learning; news report: \url{https://agn.gt/comunidad-linguistica-poqomchi-lanza-diccionario-pedagogico/}. The variety covered is not stated in the report.
\end{itemize}

\section{Linguistic and Reference Materials}
\begin{itemize}
\item Native-languages.org overview page on Pocomchi, with links to further references: \url{https://native-languages.org/pocomchi.htm}
\item Endangered Languages Project entry (variety not verified): \url{https://endangeredlanguages.com/node/112918}
\end{itemize}
No dedicated grammar was found in the searches; academic databases should be checked separately.

\section{Religious and Community Resources}
\begin{itemize}
\item Poqomchi' Bible (New Testament, 2018 edition in the reformed orthography; Wycliffe-related translation) at Digital Bible Society; the entry is tagged with a Poqomchi' code and the variety should be verified: \url{https://dbs.org/bibles/POHBSG}
\end{itemize}

\section{Digital and Audio Resources}
\begin{itemize}
\item Poqomchi' Tusq'orik, a talking dictionary with audio (Living Dictionaries); about page: \url{https://livingdictionaries.app/poqomchi/about}
\end{itemize}

\section{Curriculum Development Context}
Very little classroom material in Poqomchi' surfaced. The ALMG dictionary and the talking dictionary are the most useful starting points for vocabulary, and the Bible text offers a sizeable written corpus in the reformed orthography. Which resources belong to Western rather than Eastern Poqomchi' must be checked with the community and ALMG before use.

\end{document}